%% kapa: exam=1 round=35 year=2024 period=1 form=A subject=civil_law subject_ko=민법 seq=7 number=7
%% source: Q-Net 감정평가사 기출문제 (cst00309) artlSeq=5213229 file=raw/1008/5213229/1교시 문제지 원본.pdf
%% conversion: kordoc 4.24.51 -> tools/gen1.py (auto); review=proofread-sonnet; flags=none
\begin{question}{7}
민법상 사단법인에 관한 설명으로 옳지 않은 것은? (다툼이 있으면 판례에 따름)
\begin{choices}
  \item 설립자가 법인의 해산사유를 정하는 경우에는 정관에 그 사유를 기재하여야 한다.
  \item 사원총회 결의에 의한 정관의 해석은 정관의 규범적 의미와 다르더라도 법인의 구성원을 구속하는 효력이 있다.
  \item 사원의 지위는 정관에 달리 정함이 없으면 양도할 수 없다.
  \item 정관에 이사의 해임사유에 관한 규정이 있는 경우, 법인은 특별한 사정이 없는 한 정관에서 정하지 않은 사유로 이사를 해임할 수 없다.
  \item 법원의 직무집행정지 가처분결정에 의해 권한이 정지된 대표이사가 그 정지기간 중 체결한 계약은 그 후 가처분신청이 취하되었더라도 무효이다.
\end{choices}
\end{question}
\answer{2}
